                               IMC2007, Blagoevgrad, Bulgaria
                                          Day 1, August 5, 2007


Problem 1. Let f be a polynomial of degree 2 with integer coefficients. Suppose that f (k) is divisible
by 5 for every integer k. Prove that all coefficients of f are divisible by 5.
Solution 1. Let f (x) = ax2 + bx + c. Substituting x = 0, x = 1 and x = −1, we obtain that 5|f (0) = c,
5|f (1) = (a + b + c) and 5|f (−1) = (a − b + c). Then 5|f (1) + f (−1) − 2f (0) = 2a and 5|f (1) − f (−1) = 2b.
Therefore 5 divides 2a, 2b and c and the statement follows.
Solution 2. Consider f (x) as a polynomial over the 5-element field (i.e. modulo 5). The polynomial has
5 roots while its degree is at most 2. Therefore f ≡ 0 (mod 5) and all of its coefficients are divisible by 5.

Problem 2. Let n ≥ 2 be an integer. What is the minimal and maximal possible rank of an n × n matrix
whose n2 entries are precisely the numbers 1, 2, . . . , n2 ?
Solution. The minimal rank is 2 and the maximal rank is n. To prove this, we have to show that the rank
can be 2 and n but it cannot be 1.
   (i) The rank is at least 2. Consider an arbitrary matrix A = [aij ] with entries 1, 2, . . . , n2 in some
order. Since permuting rows or columns of a matrix does not change its rank, we can assume that
1 = a11 < a21 < · · · < an1 and a11 < a12 <
                                            · · · < a1n . Hence an1 ≥ n and a1n ≥ n and at
                                                                                           least one of these
                                   a11 a1n                                         a   a
inequalities is strict. Then det             < 1 · n2 − n · n = 0 so rk(A) ≥ rk 11 1n ≥ 2.
                                  an1 ann                                         an1 ann
   (ii) The rank can be 2. Let
                                                                              
                                               1              2      ... n
                                         n+1                n+2     . . . 2n
                                    T =        ..            ..            .. 
                                                                              
                                                                     ..
                                                .             .         .   . 
                                          n − n + 1 n − n + 2 . . . n2
                                           2               2


The ith row is (1, 2, . . . , n) + n(i − 1) · (1, 1, . . . , 1) so each row is in the two-dimensional subspace generated
by the vectors (1, 2, . . . , n) and (1, 1, . . . , 1). We already proved that the rank is at least 2, so rk(T ) = 2.
   (iii) The rank can be n, i.e. the matrix can be nonsingular. Put odd numbers into the diagonal,
only even numbers above the diagonal and arrange the entries under the diagonal arbitrarily. Then the
determinant of the matrix is odd, so the rank is complete.

Problem 3. Call a polynomial P (x1 , . . . , xk ) good if there exist 2 × 2 real matrices A1 , . . . , Ak such that
                                                               k
                                                                       !
                                                              X
                                 P (x1 , . . . , xk ) = det       xi Ai .
                                                                  i=1

Find all values of k for which all homogeneous polynomials with k variables of degree 2 are good.
    (A polynomial is homogeneous if each term has the same total degree.)
Solution. The possible values for k are 1 and 2.            
    If k = 1 then P (x) = αx2 and we can choose A1 = 10 α0 .
                                                                                                           
                                   2     2                                                1 0             0 β
    If k = 2 then P (x, y) = αx + βy + γxy and we can choose matrices A1 = 0 α and A2 = −1 γ .
                                                                                k
                                                                                   x2i is not good. Suppose that
                                                                                P
    Now let k ≥ 3. We show that the polynomial P (x1 , . . . , xk ) =
                           k                                                 i=0
                            P
P (x1 , . . . , xk ) = det    xi Ai . Since the first columns of A1 , . . . , Ak are linearly dependent, the first
                         i=0


                                                           1
column of some non-trivial linear combination y1 A1 + . . . + yk Ak is zero. Then det(y1A1 + . . . + yk Ak ) = 0
but P (y1, . . . , yk ) 6= 0, a contradiction.

Problem 4. Let G be a finite group. For arbitrary sets U, V, W ⊂ G, denote by NU V W the number of
triples (x, y, z) ∈ U × V × W for which xyz is the unity.
    Suppose that G is partitioned into three sets A, B and C (i.e. sets A, B, C are pairwise disjoint and
G = A ∪ B ∪ C). Prove that NABC = NCBA .
Solution. We start with three preliminary observations.
    Let U, V be two arbitrary subsets of G. For each x ∈ U and y ∈ V there is a unique z ∈ G for which
xyz = e. Therefore,
                                       NU V G = |U × V | = |U| · |V |.                                 (1)
Second, the equation xyz = e is equivalent to yzx = e and zxy = e. For arbitrary sets U, V, W ⊂ G, this
implies

{(x, y, z) ∈ U ×V ×W : xyz = e} = {(x, y, z) ∈ U ×V ×W : yzx = e} = {(x, y, z) ∈ U ×V ×W : zxy = e}

and therefore
                                           NU V W = NV W U = NW U V .                                       (2)
Third, if U, V ⊂ G and W1 , W2 , W3 are disjoint sets and W = W1 ∪ W2 ∪ W3 then, for arbitrary U, V ⊂ G,

            {(x, y, z) ∈ U × V × W : xyz = e} = {(x, y, z) ∈ U × V × W1 : xyz = e}∪

                     ∪{(x, y, z) ∈ U × V × W2 : xyz = e} ∪ {(x, y, z) ∈ U × V × W3 : xyz = e}
so
                                    NU V W = NU V W1 + NU V W2 + NU V W3 .                                  (3)
     Applying these observations, the statement follows as

                      NABC = NABG − NABA − NABB = |A| · |B| − NBAA − NBAB =

                                = NBAG − NBAA − NBAB = NBAC = NCBA .

Problem 5.     Let n be a positive integer and a1 , . . . , an be arbitrary integers. Suppose that a function
                    n
                    X
f : Z → R satisfies   f (k + ai ℓ) = 0 whenever k and ℓ are integers and ℓ 6= 0. Prove that f = 0.
                     i=1
Solution. Let us define a subset I of the polynomial ring R[X] as follows:
                       n          m
                                  X              m
                                                 X                                               o
                                           j
                  I = P (X) =           bj X :         bj f (k + jℓ) = 0 for all k, ℓ ∈ Z, ℓ 6= 0 .
                                  j=0            j=0

This is a subspace of the real vector space R[X]. Furthermore, P (X)P∈ I implies X · P (X) ∈ I. Hence,
                                                                         n     ai
I is an ideal, and it is non-zero, because the polynomial R(X) =         i=1 X    belongs to I. Thus, I is
generated (as an ideal) by some non-zero polynomial Q.
   If Q is constant then the definition of I implies f = 0, so we can assume that Q has a complex zero c.
Again, by the definition of I, the polynomial Q(X m ) belongs to I for every natural number m ≥ 1; hence
Q(X) divides Q(X m ). This shows that all the complex numbers

                                                   c, c2 , c3 , c4 , . . .

are roots of Q. Since Q can have only finitely many roots, we must have cN = 1 for some N ≥ 1; in
particular, Q(1) = 0, which implies P (1) = 0 for all P ∈ I. This contradicts the fact that R(X) =
P n     ai
  i=1 X    ∈ I, and we are done.

                                                             2
Problem 6. How many nonzero coefficients can a polynomial P (z) have if its coefficients are integers
and |P (z)| ≤ 2 for any complex number z of unit length?
Solution. We show that the number of nonzero coefficients can be 0, 1 and 2. These values are possible,
for example the polynomials P0 (z) = 0, P1 (z) = 1 and P2 (z) = 1 + z satisfy the conditions and they have
0, 1 and 2 nonzero terms, respectively.
    Now consider an arbitrary polynomial P (z) = a0 + a1 z + . . .+ an z n satisfying the conditions and assume
that it has at least two nonzero coefficients. Dividing the polynomial by a power of z and optionally
replacing p(z) by −p(z), we can achieve a0 > 0 such that conditions are not changed and the number of
nonzero terms is preserved. So, without loss of generality, we can assume that a0 > 0.
    Let Q(z) = a1 z + . . . + an−1 z n−1 . Our goal is to show that Q(z) = 0.
    Consider those complex numbers w0 , w1 , . . . , wn−1 on the unit circle for which an wkn = |an |; namely, let
                                      (
                                        e2kπi/n      if an > 0
                               wk =       (2k+1)πi/n
                                                                  (k = 0, 1, . . . , n).
                                        e            if an < 0

Notice that
                         n−1
                         X                n−1
                                          X                         n−1
                                                                    X              n−1
                                                                                   X
                               Q(wk ) =          Q(w0 e2kπi/n ) =         aj w0j         (e2jπi/n )k = 0.
                         k=0               k=0                      j=1            k=0

   Taking the average of polynomial P (z) at the points wk , we obtain
                              n−1             n−1
                            1X              1X
                                                  a0 + Q(wk ) + an wkn = a0 + |an |
                                                                      
                                  P (wk ) =
                            n k=0           n k=0

and
                                     n−1                    n−1
                                  1X              1X
                               2≥       P (wk ) ≥       P (wk ) = a0 + |an | ≥ 2.
                                  n k=0           n k=0

This obviously implies a0 = |an | = 1 and P (wk ) = 2 + Q(wk ) = 2 for all k. Therefore, all values of
Q(wk ) must lie on the circle |2 + z| = 2, while their sum is 0. This is possible only if Q(wk ) = 0 for all k.
Then polynomial Q(z) has at least n distinct roots while its degree is at most n − 1. So Q(z) = 0 and
P (z) = a0 + an z n has only two nonzero coefficients.
Remark. From Parseval’s formula (i.e. integrating |P (z)|2 = P (z)P (z) on the unit circle) it can be
obtained that                                       Z 2π                Z 2π
                              2               2   1          it 2     1
                         |a0 | + . . . + |an | =         P (e ) dt ≤         4 dt = 4.                      (4)
                                                 2π 0                2π 0
Hence, there cannot be more than four nonzero coefficients, and if there are more than one nonzero term,
then their coefficients are ±1.
   It is also easy to see that equality in (4) cannot hold two or more nonzero coefficients, so it is sufficient
to consider only polynomials of the form 1 ± xm ± xn . However, we do not know (yet :-)) any simpler
argument for these cases than the proof above.




                                                             3
